% F035 指数幂的运算法则｜01 核心结论
\begin{statementbox}
    \textbf{\textcolor{CStatement}{先看一个问题}}\par
    $2^3\times2^4$ 要不要把 $7$ 个 $2$ 全部乘出来？不用。把两组相同的因子接起来，就有 $2^{3+4}=2^7$。

    \medskip
    \begin{center}
    \begin{tikzpicture}[font=\small,>=stealth]
        \node[draw=green!45!black,fill=green!8,rounded corners=3pt,inner sep=6pt] (l) at (0,0) {$2\cdot2\cdot2$};
        \node[draw=green!45!black,fill=green!8,rounded corners=3pt,inner sep=6pt] (r) at (4.0,0) {$2\cdot2\cdot2\cdot2$};
        \node at (2.0,0) {$\times$};
        \node[font=\footnotesize] at (0,-0.65) {3 个 $2$};
        \node[font=\footnotesize] at (4.0,-0.65) {4 个 $2$};
        \node[draw=green!55!black,fill=green!13,rounded corners=3pt,inner sep=6pt] (ans) at (2.0,-1.7) {$2^7\quad\text{共 7 个 }2$};
        \draw[->,green!50!black,thick] (2.0,-0.9) -- (ans.north);
    \end{tikzpicture}

    \end{center}

    \textbf{\textcolor{CStatement}{三条一定要会的规律}}\par
    当 $a,b>0$，$m,n$ 为任意实数时：
    \begin{description}[style=nextline,leftmargin=2.5em,itemsep=0.45em]
        \item[\textbf{一、 同底数幂相乘：指数相加。}]
        \[a^m\cdot a^n=a^{m+n}.\]
        \textbf{看结构：} $a^m$ 和 $a^n$ 并排相乘，合并成一组。
        \item[\textbf{二、 幂的乘方：指数相乘。}]
        \[(a^m)^n=a^{mn}.\]
        \textbf{看结构：} 括号外还有一次乘方，相当于每组 $m$ 个，共有 $n$ 组。
        \item[\textbf{三、 积的乘方：每个因子都要乘方。}]
        \[(ab)^m=a^mb^m.\]
        \textbf{看结构：} $ab$ 整体乘方，$a$ 和 $b$ 都不能漏。
    \end{description}

    \textbf{\textcolor{CStatement}{使用前先看条件}}\par
    \begin{itemize}[leftmargin=*,itemsep=0.25em,topsep=0.3em]
        \item \textbf{任意实数指数：} 用 $a,b>0$ 这一组安全条件。
        \item \textbf{整数指数：} 若底数非零，正数和负数都可以套用；正整数次幂的底数还可以是 $0$。
        \item \textbf{出现零次幂或负整数次幂：} 底数不能是 $0$，因为不能除以 $0$。
    \end{itemize}
\end{statementbox}
